%% parts/p07/normalization-and-curry-howard-supplement.tex
% Necessity: the exact inherited normalization chapter restores its own STLC
% and Curry--Howard text but does not contain the separately attributed
% categorical dictionary and A-072/A-151/A-152 visual apparatus.
% Definition: this separately authored derivative begins from historical
% lines 150--358 and now supplies the exact product/sum/unit/empty typed
% extension required by its bicartesian closed categorical construction.
% Standing: it is an editorial reader supplement, not inherited source, an
% admitted theorem family, or a proof-status upgrade.
% Nonredundancy: the exact source controls its theorem and worked-example
% surfaces; this block adds only the categorical constructions and diagrams.
% Use: the Part VII driver routes it after all five exact normalization
% wrappers; it is not represented as a verbatim restored slice.

% BEGIN HISTORICAL SLICE 150-358
\section{Curry--Howard and categorical semantics}

\begin{hbdefinition}{def:p07-curry-howard}{Curry--Howard setting and categorical standing}
The \emph{Curry--Howard correspondence} used here is the structural
translation between a fixed intuitionistic natural-deduction calculus and a
fixed typed term calculus: formulas map to types, assumption occurrences to
typed variables, derivations to typing derivations, discharge to abstraction,
elimination to application, and normalization to computation.  It compares
specified formal objects; it does not identify every proof with a program
without a term assignment and conversion relation.

For the third column of Table~\ref{tab:p07-curry-howard}, a
\emph{bicartesian closed category} means a category with finite products
(hence a terminal object), finite coproducts (hence an initial object), and,
for each object \(X\), a right adjoint to \(-\times X\) (hence
exponentials).  This hypothesis is necessary because products, sums, truth,
falsity, and implication have no uniform readings in an arbitrary category.

The arrow-only calculus on \texttt{sheet:stlc} cannot generate that structure.
For the categorical column only, use its following explicitly defined
extension.  Types are generated by
\[
 \alpha::=\mathsf b\mid\mathsf 1_{\mathrm{ty}}\mid\mathsf 0_{\mathrm{ty}}
          \mid\alpha\mathbin{\times_{\mathrm{ty}}}\beta
          \mid\alpha\mathbin{+_{\mathrm{ty}}}\beta
          \mid\alpha\mathbin{\to_{\mathrm{ty}}}\beta.
\]
Besides variables, abstraction, and application, terms contain
\(\star:\mathsf 1_{\mathrm{ty}}\), pairs \(\langle M,N\rangle\), projections
\(\pi_1M,\pi_2M\), injections \(\mathsf{inl}(M),\mathsf{inr}(N)\),
\(\mathsf{case}(L;x.M;y.N)\), and \(\mathsf{abort}_{\alpha}(M)\).
Their typing rules are exactly the rows of the following table; every listed
premise is required.
\begin{center}
\begingroup\small
\begin{tabularx}{\textwidth}{@{}l X X@{}}
\toprule
Constructor & Premises & Conclusion \\
\midrule
pair & \(\Theta\vdash M:\alpha\), \(\Theta\vdash N:\beta\) &
\(\Theta\vdash\langle M,N\rangle:
  \alpha\mathbin{\times_{\mathrm{ty}}}\beta\) \\
projection \(i\) & \(\Theta\vdash P:
  \alpha\mathbin{\times_{\mathrm{ty}}}\beta\) &
\(\Theta\vdash\pi_iP:\alpha_i\), where
\(\alpha_1=\alpha\), \(\alpha_2=\beta\) \\
left injection & \(\Theta\vdash M:\alpha\) &
\(\Theta\vdash\mathsf{inl}(M):\alpha\mathbin{+_{\mathrm{ty}}}\beta\) \\
right injection & \(\Theta\vdash N:\beta\) &
\(\Theta\vdash\mathsf{inr}(N):\alpha\mathbin{+_{\mathrm{ty}}}\beta\) \\
case & \(\Theta\vdash L:\alpha\mathbin{+_{\mathrm{ty}}}\beta\),
\(\Theta,x:\alpha\vdash M:\gamma\),
\(\Theta,y:\beta\vdash N:\gamma\) &
\(\Theta\vdash\mathsf{case}(L;x.M;y.N):\gamma\) \\
empty elimination & \(\Theta\vdash M:\mathsf 0_{\mathrm{ty}}\) &
\(\Theta\vdash\mathsf{abort}_{\alpha}(M):\alpha\) \\
unit & none & \(\Theta\vdash\star:\mathsf 1_{\mathrm{ty}}\) \\
\bottomrule
\end{tabularx}
\endgroup
\end{center}
The arrow rules are those of \texttt{sheet:stlc}.  Equality is the least typed
congruence, stable under capture-avoiding substitution, containing arrow
beta--eta, the two projection beta equations and
\(\langle\pi_1P,\pi_2P\rangle=P\), the two case beta equations and
\[
 \mathsf{case}(L;x.\mathsf{inl}(x);y.\mathsf{inr}(y))=L,
\]
together with \(M=\star\) at type \(\mathsf 1_{\mathrm{ty}}\) and
\(M=\mathsf{abort}_{\alpha}(z)\) for every
\(z:\mathsf 0_{\mathrm{ty}}\vdash M:\alpha\).  This is a formal typed calculus used by the
categorical construction, not an alteration of the registered arrow-only
sheet.  It is necessary and nonredundant precisely because the four added type
formers supply the product, coproduct, terminal, and initial universal
properties used below.  The displayed least congruence is the
\emph{beta--eta equivalence} used in the table.  It is a metalinguistic
relation on well-typed terms, not an additional directed reduction rule; its
eta and uniqueness equations distinguish the categorical quotient from the
beta-only operational semantics of \texttt{sheet:stlc}.

The resulting syntactic category has these types as objects.  A morphism from
\(\alpha\) to \(\beta\) is an equality class of judgments
\(x{:}\alpha\vdash M{:}\beta\); the variable judgment is the identity and
capture-avoiding substitution is composition.  The displayed beta--eta and
uniqueness equations prove that this category is bicartesian closed.
Conversely, the \emph{internal-language construction} used in the figure
represents each object of a chosen bicartesian closed category by a type and
each morphism from \(X\) to \(Y\) by a typed unary operation named \(f\), with
the chosen products, coproducts, and exponentials supplying the corresponding
constructors.  Both
are metalinguistic constructions relative to fixed data: the syntactic
category sends syntax to a category, whereas the internal language sends a
chosen category to syntax.  The distinction is used on the two directions of
the categorical edge in the figure.  These definitions also distinguish
term assignment from categorical semantics and are used by the table and
Figure~\ref{fig:p07-curry-howard-triangle}.
\index{Curry--Howard correspondence}
\index{bicartesian closed category}
\end{hbdefinition}

\begingroup\small
\begin{longtable}{@{}p{.24\textwidth}p{.29\textwidth}p{.38\textwidth}@{}}
\caption{Logic, type theory, and category dictionary.}
\label{tab:p07-curry-howard}\\
\toprule
Intuitionistic logic & Typed lambda calculus & Chosen bicartesian closed
category \\
\midrule
\endfirsthead
\toprule
Intuitionistic logic & Typed lambda calculus & Chosen bicartesian closed
category \\
\midrule
\endhead
Formula & Type & Object \\
Assumption context & Typed context, represented by a product of types &
Context object, represented by the corresponding categorical product \\
Derivation of \(\varphi\) from \(\Gamma\) & Term \(M\) of type \(\alpha\)
in a typed context \(\Theta_{\mathrm{ty}}\) & Morphism from the context object
to the object interpreting \(\varphi\) \\
Implication \(\varphi\to\psi\) & Function type
\(\alpha\mathbin{\to_{\mathrm{ty}}}\beta\) & Exponential \(Y^X\), where
\(X,Y\) interpret \(\varphi,\psi\) \\
Conjunction & Product type & Categorical product \\
Truth & Unit type & Terminal object \\
Disjunction & Sum type & Categorical coproduct \\
Falsity & Empty type & Initial object \\
Proof substitution / Cut & Capture-avoiding term substitution & Morphism
composition \\
Beta--eta conversion & Term equality after the declared conversions &
Equality of denoted morphisms \\
\bottomrule
\end{longtable}
\endgroup

The product, sum, unit, and empty-type rows use the exact extension in
Definition~\ref{def:p07-curry-howard}; they do not add constructors or
reduction rules to \texttt{sheet:stlc}.  Categorical details are supported by
\cite{GirardLafontTaylor1989,LambekScott1986} and receive full treatment in
the later categorical-logic part.

\begin{figure}[!htbp]
\centering
\resizebox{.9\linewidth}{!}{%
\begin{tikzpicture}[font=\small,>=Latex]
  \node[draw,rounded corners,align=center] (logic) at (0,2.8)
    {intuitionistic logic\\formulas and derivations};
  \node[draw,rounded corners,align=center] (types) at (-4.2,0)
    {typed lambda calculus\\types and terms modulo conversions};
  \node[draw,rounded corners,align=center] (cats) at (4.2,0)
    {bicartesian closed category\\objects and morphisms};
  \draw[<->] (logic) -- node[above left,align=center]
    {Curry--Howard\\term assignment} (types);
  \draw[<->] (types) -- node[below,align=center]
    {categorical semantics} (cats);
  \draw[<->] (cats) -- node[above right,align=center]
    {internal language /\\syntactic category} (logic);
\end{tikzpicture}
}
\caption{The Curry--Howard--categorical triangle.  Each edge
is a specified construction, not an equality sign; the categorical vertex
has the hypotheses of Definition~\ref{def:p07-curry-howard}.}
\label{fig:p07-curry-howard-triangle}
\end{figure}

\clearpage
\begin{landscape}
\begin{figure}[p]
\centering
\begin{minipage}[t]{.49\linewidth}
\centering\scriptsize
\textbf{STLC typing instance}\\[4pt]
\begin{tikzpicture}[font=\scriptsize,>=Latex,every node/.style={align=center}]
  \node (f) at (-2.1,4.2)
    {\(\Theta_{\mathrm{ty}}\derives_{\mathsf{STLC}}f{:}
      \beta\mathbin{\to_{\mathrm{ty}}}\gamma\)};
  \node (g) at (.2,4.2)
    {\(\Theta_{\mathrm{ty}}\derives_{\mathsf{STLC}}g{:}
      \alpha\mathbin{\to_{\mathrm{ty}}}\beta\)};
  \node (x) at (2.2,4.2)
    {\(\Theta_{\mathrm{ty}}\derives_{\mathsf{STLC}}x{:}\alpha\)};
  \node (gx) at (1.2,3.25)
    {\(\Theta_{\mathrm{ty}}\derives_{\mathsf{STLC}}gx{:}\beta\)};
  \node (fgx) at (0,2.3)
    {\(\Theta_{\mathrm{ty}}\derives_{\mathsf{STLC}}f(gx){:}\gamma\)};
  \node (lx) at (0,1.45)
    {\(f{:}\beta\mathbin{\to_{\mathrm{ty}}}\gamma,\,
      g{:}\alpha\mathbin{\to_{\mathrm{ty}}}\beta
      \derives_{\mathsf{STLC}}\lambda x{:}\alpha.f(gx){:}
      \alpha\mathbin{\to_{\mathrm{ty}}}\gamma\)};
  \node (lg) at (0,.65)
    {\(f{:}\beta\mathbin{\to_{\mathrm{ty}}}\gamma
      \derives_{\mathsf{STLC}}
      \lambda g{:}\alpha\mathbin{\to_{\mathrm{ty}}}\beta.
      \lambda x{:}\alpha.f(gx){:}
      (\alpha\mathbin{\to_{\mathrm{ty}}}\beta)
      \mathbin{\to_{\mathrm{ty}}}
      (\alpha\mathbin{\to_{\mathrm{ty}}}\gamma)\)};
  \node (lf) at (0,-.25)
    {\(\varnothing\derives_{\mathsf{STLC}}
      \lambda f{:}\beta\mathbin{\to_{\mathrm{ty}}}\gamma.
      \lambda g{:}\alpha\mathbin{\to_{\mathrm{ty}}}\beta.
      \lambda x{:}\alpha.f(gx){:}
      (\beta\mathbin{\to_{\mathrm{ty}}}\gamma)
      \mathbin{\to_{\mathrm{ty}}}
      ((\alpha\mathbin{\to_{\mathrm{ty}}}\beta)
      \mathbin{\to_{\mathrm{ty}}}
      (\alpha\mathbin{\to_{\mathrm{ty}}}\gamma))\)};
  \draw[->] (g) -- (gx);
  \draw[->] (x) -- node[right]
    {\hyperref[rule:stlc-app]{\(\to_{\mathrm{ty}}E\)}} (gx);
  \draw[->] (f) -- (fgx);
  \draw[->] (gx) -- node[right]
    {\hyperref[rule:stlc-app]{\(\to_{\mathrm{ty}}E\)}} (fgx);
  \draw[->] (fgx) -- node[right]
    {\hyperref[rule:stlc-abs]{\(\to_{\mathrm{ty}}I\)}} (lx);
  \draw[->] (lx) -- node[right]
    {\hyperref[rule:stlc-abs]{\(\to_{\mathrm{ty}}I\)}} (lg);
  \draw[->] (lg) -- node[right]
    {\hyperref[rule:stlc-abs]{\(\to_{\mathrm{ty}}I\)}} (lf);
\end{tikzpicture}
\end{minipage}\hfill
\begin{minipage}[t]{.49\linewidth}
\centering\scriptsize
\textbf{Intuitionistic ND instance}\\[4pt]
\begin{tikzpicture}[font=\scriptsize,>=Latex,every node/.style={align=center}]
  \node (f) at (-2.1,4.2) {\([\psi\to\chi]^1\)};
  \node (g) at (.2,4.2) {\([\varphi\to\psi]^2\)};
  \node (x) at (2.2,4.2) {\([\varphi]^3\)};
  \node (b) at (1.2,3.25) {\(\psi\)};
  \node (c) at (0,2.3) {\(\chi\)};
  \node (ac) at (0,1.45) {\(\varphi\to\chi\)};
  \node (gb) at (0,.65)
    {\((\varphi\to\psi)\to(\varphi\to\chi)\)};
  \node (all) at (0,-.25)
    {\((\psi\to\chi)\to
      ((\varphi\to\psi)\to(\varphi\to\chi))\)};
  \draw[->] (g) -- (b);
  \draw[->] (x) -- node[right]
    {\hyperref[rule:nj-imp-e]{\(\to E\)}} (b);
  \draw[->] (f) -- (c);
  \draw[->] (b) -- node[right]
    {\hyperref[rule:nj-imp-e]{\(\to E\)}} (c);
  \draw[->] (c) -- node[right]
    {\hyperref[rule:nj-imp-i]{\(\to I,3\)}} (ac);
  \draw[->] (ac) -- node[right]
    {\hyperref[rule:nj-imp-i]{\(\to I,2\)}} (gb);
  \draw[->] (gb) -- node[right]
    {\hyperref[rule:nj-imp-i]{\(\to I,1\)}} (all);
\end{tikzpicture}
\end{minipage}
\caption{Composition as one complete STLC typing instance beside its
intuitionistic natural-deduction instance.  Generic rules
remain on \texttt{sheet:stlc} and \texttt{sheet:nj}; arrow labels point to
their registered rows.  Here
\(\Theta_{\mathrm{ty}}=
f{:}\beta\mathbin{\to_{\mathrm{ty}}}\gamma,\,
g{:}\alpha\mathbin{\to_{\mathrm{ty}}}\beta,\,x{:}\alpha\).}
\label{fig:p07-stlc-nd-derivation}
\end{figure}
\end{landscape}
\clearpage

% END HISTORICAL SLICE 150-358
